\begin{abstract}
Issuing reliable forecast bands for astronomical light curves requires quantifying future brightness uncertainty from early detections before an object's astrophysical class is spectroscopically confirmed.
Deployed pipelines such as FLARE route conformal multipliers by the \emph{predicted} class and therefore under-cover rare transients when early classifications err---capturing only \ZTFflareSLSNsixtyeight{} of future superluminous supernova (SLSN) measurements at the nominal 68\% level from a five-day prefix---whereas existing latent-group valid methods preserve coverage only by discarding posterior structure and inflating band width.
Guided by the population form of the shortest scale-weighted band satisfying per-class coverage constraints, we introduce \gamut{}, a conformal calibration framework that constructs prediction regions as superlevel sets of a dual-weighted posterior predictive density mixture.
Specifically, \gamut{} couples a shared hypothesis-conditional quantile forecaster and leave-one-object-out residual densities with temperature-scaled membership posteriors and a multiplicative fixed-point solver that equalises per-class conformal risk-control thresholds; evaluated under cross-conformal data reuse it approximates the oracle band empirically, while under a strict split with disjoint threshold certification it delivers finite-sample coverage for every class--horizon cell with sufficient certification objects.
Across six time-series benchmarks, 13 baselines, a strict three-way split, and two pre-registered held-out LSST sets, \gamut{} raises day-5 ZTF SLSN coverage from \ZTFflareSLSNsixtyeight{} to \ZTFgamutSLSNsixtyeight{} at \ZTFfiveWidthRatioSixtyeight{} times FLARE's width, leaves no class significantly under-covered in \MainValidGamutSixtyeight{} (68\%) and \MainValidGamutNinetyfive{} (95\%) of main settings, and shortens bands by \WidthGainValidSixtyeight{} (68\%) and \WidthGainValidNinetyfive{} (95\%) on average relative to always-valid baselines, widening only when 95\% risk budgets meet very small classes.
In a companion study predicting whether one further detection narrows FLARE's 90\% classification set to a singleton, combining learned next-detection laws with Venn--Abers recalibration lowers test Brier score by \CloFiveBrierGainVA{} (day 5) and \CloTenBrierGainVA{} (day 10), driven primarily by replacing the fixed-scatter Gaussian simulation rather than class mixing.
\end{abstract}
